\section*{الفصل \ref{ch:g12:special-relativity} --- النسبية الخاصة: تمدد الزمن}

\begin{solution}{exo:g12:special-relativity:1}
$\gamma = 1/\sqrt{1 - (v/c)^2}$: $1.01$ ($0.1c$)؛ $1.25$ ($0.6c$)؛
$1.67$ ($0.8c$)؛ $10.0$ ($0.995c$).
\end{solution}

\begin{solution}{exo:g12:special-relativity:2}
$\gamma = 1.25$، ومن ثم تقيس الأرض $1.25 \times 10.0 =
\qty{12.5}{s}$. والمقدار \qty{10.0}{s} في السفينة هو \omterm{def:g12:special-relativity:proper-time}{الزمن الذاتي}:
فالميقاتية (وهي ساعة واحدة) حاضرة عند النبضتين.
\end{solution}

\begin{solution}{exo:g12:special-relativity:3}
$v = c\sqrt{1 - 1/\gamma^2}$: يعطي $\gamma = 2$ المقدارَ $v = 0.866c$؛
ويعطي $\gamma = 10$ المقدارَ $v = 0.995c$.
\end{solution}

\begin{solution}{exo:g12:special-relativity:4}
$v = \qty{36.1}{m/s}$، $v/c = \num{1.2e-7}$؛
$\gamma - 1 \approx (\num{1.2e-7})^2/2 \approx \num{7.2e-15}$. وتُفقد \omterm{def:g10:orders-of-magnitude:unit}{ثانية}
واحدة بعد $1/\num{7.2e-15} \approx \qty{1.4e14}{s}$ --- أي نحو
أربعة ملايين سنة من القيادة.
\end{solution}

\begin{solution}{exo:g12:special-relativity:5}
يركب الراكب نحو الوميض الأمامي، ومن ثم يبلغه أولًا؛
وبالمسلَّمة الثانية يسافر الوميضان \omterm{def:g12:kinematics-2d:velocity}{بالسرعة} $c$ نفسها على مسافتَي نصف القطار
المتساويتين \emph{في مرجعه}، ومن ثم يعني الوصول الأبكر
ضربةً أبكر. فالتزامن يتوقف على \omterm{def:g10:relative-motion:frame}{المرجع}: وكل حكم
صحيح في مرجعه.
\end{solution}

\begin{solution}{exo:g12:special-relativity:6}
$\gamma = 1/\sqrt{1 - 0.99^2} = 7.09$؛ والعمر في المختبر
$7.09 \times \num{2.6e-8} = \qty{1.8e-7}{s}$. \omterm{def:g12:projectile-motion:range}{والمدى} الكلاسيكي
$0.99 \times \num{3.00e8} \times \num{2.6e-8} \approx \qty{7.7}{m}$؛
والفعلي $7.09 \times 7.7 \approx \qty{55}{m}$.
\end{solution}

\begin{solution}{exo:g12:special-relativity:7}
$\gamma = 1.5$ عند $v = c\sqrt{1 - 1/2.25} \approx 0.75c$؛
$\gamma = 3$ عند $\approx 0.94c$. أما المقارب: فالعامل $\gamma$ (ومعه
كلفة \omterm{def:g10:energy-conservation:energy}{الطاقة} لمزيد من \omterm{def:g12:kinematics-2d:velocity}{السرعة}) يتباعد عند $v \to c$ --- فلا يبلغ أي
جهد منتهٍ بجسم ذي كتلة سرعةَ الضوء.
\end{solution}

\begin{solution}{exo:g12:special-relativity:8}
$\Delta t_0 = 2L/c = 3.0/\num{3.00e8} = \qty{1.0e-8}{s}$؛ وعند $0.8c$،
$\gamma = 5/3$، ومنه $\Delta t = \qty{1.7e-8}{s}$.
\end{solution}

\begin{solution}{exo:g12:special-relativity:9}
\omterm{def:g12:special-relativity:proper-time}{الزمن الذاتي} يخصّ الساعة الحاضرة عند الحدثين: (أ) اللابس؛
(ب) الميون؛ (ج) الطيار --- فالسفينة وحدها حاضرة عند
المغادرة وعند الوصول.
\end{solution}

\begin{solution}{exo:g12:special-relativity:10}
$\gamma - 1 \approx 7700^2/(2 \times \num{9.0e16}) \approx
\num{3.3e-10}$؛ وعلى مدى \qty{1.58e7}{s} تتأخر رائدة الفضاء
$\num{3.3e-10} \times \num{1.58e7} \approx \qty{5.2e-3}{s}$ --- أي أصغر بنحو
خمس ميلي ثوانٍ.
\end{solution}

\begin{solution}{exo:g12:special-relativity:11}
$E = \num{1.0e-3} \times (\num{3.00e8})^2 = \qty{9.0e13}{J}$. و
محطة قدرها \qty{1.0}{GW} توصّل $\num{1.0e9} \times 86400 \approx
\qty{8.6e13}{J}$ في اليوم --- فالغرام الواحد هو إنتاج يوم واحد منها، و
$\num{9.0e13}/\num{8.5e5} \approx \num{e8}$: أي مئة مليون
سيارة مسرعة.
\end{solution}

\begin{solution}{exo:g12:special-relativity:12}
ربّع واقلب: $\gamma^2(1 - v^2/c^2) = 1$، ومنه
$v^2/c^2 = 1 - 1/\gamma^2$، ومن ثم $v = c\sqrt{1 - 1/\gamma^2}$. ومن أجل
$\gamma = 30$: $v = c\sqrt{1 - 1/900} = 0.99944c$.
\end{solution}

\begin{solution}{exo:g12:special-relativity:13}
السُّمك المتقلص $L = \qty{15}{km}/30 = \qty{500}{m}$؛ وزمن المرور
$500/(0.99944 \times \num{3.00e8}) \approx \qty{1.7e-6}{s} <
\qty{2.2}{\micro s}$. فالمرجعان يتفقان على أن الميون يبلغ الأرض ---
هذا يلوم عمرًا ممطوطًا، وذاك جوًّا متقلصًا.
\end{solution}

\begin{solution}{exo:g12:special-relativity:14}
في المختبر: $3200/\num{3.00e8} \approx \qty{1.1e-5}{s}$ (\omterm{def:g12:kinematics-2d:velocity}{والسرعة}
$\approx c$). وفي \omterm{def:g10:relative-motion:frame}{مرجع} الإلكترون: الأنبوب متقلص إلى
$3200/\num{1.0e5} = \qty{3.2}{cm}$، ويُقطع في
$\num{1.1e-5}/\num{1.0e5} \approx \qty{1.1e-10}{s}$.
\end{solution}

\begin{solution}{exo:g12:special-relativity:15}
على الأرض: $4.2/0.9 \approx \qty{4.7}{years}$. ومع
$\gamma = 1/\sqrt{1 - 0.81} = 2.29$، تسجّل ساعة المسبار
$4.7/2.29 \approx \qty{2.0}{years}$.
\end{solution}

\begin{solution}{pb:g12:special-relativity:1}
\textbf{1.} $c\,\Delta t_0 = \num{3.00e8} \times \num{2.2e-6} \approx
\qty{660}{m}$.

\textbf{2.} $\num{15000}/660 \approx 23$ عمرًا.

\textbf{3.} $\mathrm{e}^{-23} \approx \num{e-10}$: أي إنه ينبغي أن يصل أقل من ميون واحد
في كل عشرة مليارات، ومع ذلك يبلغ التدفق واحدًا في \unit{cm^2} في
الدقيقة --- فالسماء تناقض الفيزياء الكلاسيكية مناقضة صريحة.

\textbf{4.} العمر على الأرض $30 \times \qty{2.2}{\micro s} =
\qty{66}{\micro s}$؛ \omterm{def:g12:projectile-motion:range}{والمدى} $\approx \num{3.00e8} \times \num{6.6e-5}
\approx \qty{20}{km} > \qty{15}{km}$: فالوصول مفسَّر.

\textbf{5.} $v = c\sqrt{1 - 1/900} = 0.99944c$.

\textbf{6.} $\Delta t_0 = 2L/c = 3.0/\num{3.00e8} = \qty{1.0e-8}{s}$.

\textbf{7.} الأفقي $v\,\Delta t/2$؛ والشاقولي $L$؛ ويقيس نصف \omterm{def:g10:relative-motion:trajectory}{المسار}
المائل $c\,\Delta t/2$ لأن المسلَّمة الثانية تثبّت
\omterm{def:g12:kinematics-2d:velocity}{سرعة} \omterm{def:g12:quantum-world:photon}{الفوتون} عند $c$ في \omterm{def:g10:relative-motion:frame}{مرجع} الأرض أيضًا.

\textbf{8.} يعطي $(c\,\Delta t/2)^2 = L^2 + (v\,\Delta t/2)^2$ المقدارَ
$\Delta t^2(c^2 - v^2) = 4L^2$، ومنه $\Delta t =
(2L/c)/\sqrt{1 - v^2/c^2} = \gamma\,\Delta t_0$.

\textbf{9.} $\gamma = 1/\sqrt{1 - 0.36} = 1.25$؛
$\Delta t = \qty{1.25e-8}{s}$.

\textbf{10.} ساعة السفينة: فساعتها الواحدة حاضرة عند الدقّتين، ومن ثم
تقرأ \omterm{def:g12:special-relativity:proper-time}{الزمن الذاتي}؛ أما الأرض فتحتاج ساعتين متزامنتين.

\textbf{11.} $\gamma = 1/\sqrt{1 - 0.64} = 1/0.6 = 5/3 \approx 1.67$.

\textbf{12.} $10.0/(5/3) = \qty{6.0}{years}$.

\textbf{13.} الذهاب $\qty{5.0}{years}$ عند $0.8c$: أي $4.0$ \omterm{def:g10:orders-of-magnitude:lightyear}{سنة ضوئية}.

\textbf{14.} $10.0 - 6.0 = \qty{4.0}{years}$: فرائدة الفضاء أصغر من توأمها بأربع
سنوات.

\textbf{15.} الوضع ليس متناظرًا: فرائدة الفضاء وحدها هي التي
تستدير --- إذ تتسارع وتغيّر مرجعها العُطالي، بينما يبقى التوأم على
الأرض في \omterm{def:g10:relative-motion:frame}{مرجع} واحد --- ومن ثم لا تسجّل الزمن الأقصر إلا ساعتها.

\textbf{16.} $v = 2\pi \times \num{2.66e7}/\num{43200} \approx
\qty{3.9e3}{m/s}$.

\textbf{17.} $\gamma - 1 \approx (\num{3.87e3})^2/(2 \times
\num{9.0e16}) \approx \num{8.3e-11}$.

\textbf{18.} $\num{8.3e-11} \times \num{86400} \approx \qty{7.2e-6}{s}
= \qty{7.2}{\micro s}$ في اليوم تأخّرًا.

\textbf{19.} $c \times \num{7.2e-6} \approx \qty{2.2}{km}$ من خطأ
الموضع في اليوم.

\textbf{20.} الانحراف الصافي $45 - 7.2 \approx \qty{38}{\micro s}$ في اليوم
تقدّمًا؛ والخطأ $c \times \num{38e-6} \approx \qty{11}{km}$ في اليوم ---
أي فاتورة النسبية التي تدفعها كل ساعة GPS مقدَّمًا، بإزاحة ضبطها
على الأرض لتدقّ صادقةً في \omterm{def:g10:universal-gravitation:satellite}{المدار}.
\end{solution}
